\documentclass[11pt]{article}
\usepackage[margin=1in]{geometry}
\usepackage{amsmath,amssymb,amsthm}
\usepackage[T1]{fontenc}
\usepackage{lmodern}
\usepackage{hyperref}

\newtheorem{theorem}{Theorem}
\newtheorem{lemma}{Lemma}
\newtheorem{proposition}{Proposition}
\newtheorem{corollary}{Corollary}

\DeclareMathOperator{\Ran}{Ran}
\DeclareMathOperator{\Tr}{Tr}
\DeclareMathOperator{\spanop}{span}

\newcommand{\Hsa}{\mathcal H_{\mathrm{sa}}}
\newcommand{\HS}{S_2}
\newcommand{\K}{\mathcal K}
\newcommand{\R}{\mathbb R}
\newcommand{\N}{\mathbb N}

\title{A positive Hilbert--Schmidt cone counterexample to a common quasi-interior lower bound}
\author{FA/Banach run \texttt{fa\_banach\_001}, agent \texttt{agent\_lane\_08}}
\date{2026-07-03}

\begin{document}
\maketitle

\section*{Source Question}

The source paper is J. Glueck and A. Mironchenko, \emph{Stability criteria for positive semigroups on ordered Banach spaces}, arXiv:2210.13566.  After Theorem 8.4, in the setting of an ordered Hilbert space $(H,H^+)$ satisfying
\[
        H^+ \subseteq (H')^+,
\]
the authors note that their hypothesis gives quasi-interior points $u,v$ whenever there is a quasi-interior point $w$ with $u,v\geq w$.  They then write that they do not know whether the converse holds: given arbitrary quasi-interior points $u,v\in H^+$, must there always be a quasi-interior point $w\in H^+$ such that $u,v\geq w$?

\section*{Result}

The answer is negative.  In fact, the following example gives two quasi-interior points with no nonzero common positive lower bound.

\begin{theorem}
There is an ordered real Hilbert space $(\Hsa,\Hsa^+)$ with self-dual cone and two quasi-interior points $A,B\in\Hsa^+$ such that
\[
        0\leq C\leq A,\ B
        \quad\Longrightarrow\quad
        C=0.
\]
Consequently there is no quasi-interior point $C$ satisfying $A,B\geq C$.
\end{theorem}

\section*{Proof Intuition}

The positive cone of self-adjoint Hilbert--Schmidt operators behaves like the infinite-dimensional positive semidefinite matrix cone.  A positive Hilbert--Schmidt operator is quasi-interior precisely when it has no kernel: its square-root range is then dense and supplies enough rank-one positive operators to densely generate the Hilbert--Schmidt space.

The obstruction is that two dense operator ranges can be transverse.  We choose an injective diagonal compact operator whose square-root range is
\[
        R=\{x=(x_n)\in\ell_2:\sum n^2 |x_n|^2<\infty\}.
\]
A Baire argument gives a unitary $U$ such that $R\cap U R=\{0\}$.  If $A$ has square-root range $R$ and $B=UAU^*$ has square-root range $UR$, Douglas' range inclusion lemma forces every positive common lower bound $C$ to have
\[
        \Ran(C^{1/2})\subseteq R\cap U R=\{0\}.
\]
Thus $C=0$, even though both $A$ and $B$ are quasi-interior.

\section*{A transverse weighted range}

We first record the elementary Hilbert-space ingredient used in the construction.

\begin{lemma}[A transverse compactly embedded dense range]\label{lem:transverse-range}
Let $E=\ell_2(\N)$ and
\[
        R=\left\{x=(x_n)\in E:\sum_{n=1}^\infty n^2 |x_n|^2<\infty\right\}.
\]
There is a unitary operator $U$ on $E$ such that
\[
        R\cap U R=\{0\}.
\]
\end{lemma}

\begin{proof}
Endow $R$ with the Hilbert norm
\[
        \|x\|_R^2=\sum_{n=1}^\infty n^2 |x_n|^2.
\]
The inclusion $R\hookrightarrow E$ is compact.  For $m\in\N$ set
\[
        K_m=\{x\in R:\|x\|_E=1,\ \|x\|_R\leq m\}.
\]
Each $K_m$ is compact in $E$, and the unit sphere of $R$ is $\bigcup_m K_m$.

Let $\mathcal U(E)$ be the unitary group with the operator-norm topology.  For $m,n\in\N$ define
\[
        \mathcal E_{m,n}
        =
        \{V\in\mathcal U(E): V K_m\cap K_n\neq\emptyset\}.
\]
The sets $\mathcal E_{m,n}$ are closed.  Indeed, if $V_j\to V$ in norm and $V_jx_j\in K_n$ for $x_j\in K_m$, compactness of $K_m$ gives a subsequence $x_j\to x\in K_m$, hence $V_jx_j\to Vx\in K_n$.

We show that each $\mathcal E_{m,n}$ has empty interior.  It suffices to prove the following small perturbation fact: if $K,L$ are compact subsets of the unit sphere of an infinite-dimensional Hilbert space, $V_0$ is unitary, and $\varepsilon>0$, then there is a unitary $V$ with $\|V-V_0\|<\varepsilon$ and $V K\cap L=\emptyset$.

Put $K_0=V_0K$.  Choose $0<\theta<\varepsilon/2$.  Since $K_0\cup L$ is compact, there is a finite-dimensional subspace $F\subset E$ such that every point of $K_0\cup L$ is within $\delta$ of $F$, where $\delta>0$ will be chosen so small that
\[
        \sin\theta(1-\delta)>2\delta.
\]
Choose an isometric copy $J:F\to F^\perp$.  Let $W$ be the unitary which rotates $F$ into $JF$ through angle $\theta$,
\[
        W f=\cos\theta\,f+\sin\theta\,Jf\qquad(f\in F),
\]
with the inverse rotation on $JF$ and the identity on $(F\oplus JF)^\perp$.
Then $\|W-I\|<\varepsilon$.  If $x\in K_0$ and $y\in L$, write $P$ for the orthogonal projection onto $F$.  The $JF$-component of $Wx$ has norm at least $\sin\theta\,\|Px\|-\|(I-P)x\|$, which is $>\delta$ by the choice of $\delta$, while the $JF$-component of $y$ has norm at most $\delta$.  Thus $Wx\neq y$.  Hence $WK_0\cap L=\emptyset$, and $V=WV_0$ has the desired properties.

Applying this perturbation fact with $K=K_m$ and $L=K_n$ proves that each $\mathcal E_{m,n}$ is nowhere dense.  Since $\mathcal U(E)$ is a complete metric space in operator norm, Baire's theorem gives a unitary
\[
        U\notin \bigcup_{m,n=1}^\infty \mathcal E_{m,n}.
\]
If $0\neq x\in R\cap UR$, then after normalizing $x$ to norm one, $x\in K_n$ for some $n$ and $x=Uy$ for a norm-one $y\in R$, with $y\in K_m$ for some $m$.  This would put $U$ in $\mathcal E_{m,n}$, a contradiction.
\end{proof}

\section*{The Hilbert--Schmidt cone}

Let $E=\ell_2(\N)$, now viewed as a real Hilbert space.  Let $\Hsa$ be the real Hilbert space of all self-adjoint Hilbert--Schmidt operators on $E$, with inner product
\[
        \langle S,T\rangle_{\Hsa}=\Tr(ST).
\]
Let $\Hsa^+$ be the cone of positive semidefinite Hilbert--Schmidt operators.

\begin{lemma}\label{lem:hs-cone}
The cone $\Hsa^+$ is self-dual in $\Hsa$.  Moreover, for $A\in\Hsa^+$, the following are equivalent:
\[
        A\text{ is a quasi-interior point of }\Hsa^+
        \quad\Longleftrightarrow\quad
        \ker A=\{0\}.
\]
\end{lemma}

\begin{proof}
Self-duality follows from rank-one tests.  If $A\in\Hsa$ satisfies $\Tr(AB)\geq0$ for every $B\in\Hsa^+$, then testing against $B=x\otimes x$ gives $\langle Ax,x\rangle\geq0$ for all $x\in E$, so $A\geq0$.  The converse implication $\Tr(AB)\geq0$ for $A,B\geq0$ follows by writing the trace as the trace of $A^{1/2}BA^{1/2}$, first for finite-rank approximants and then by approximation.

Assume first that $\ker A\neq\{0\}$ and choose $0\neq h\in\ker A$.  If $S\in[-cA,cA]$ for some $c>0$, then $cA+S\geq0$ and $cA-S\geq0$.  Since both positive operators have quadratic form value $0$ at $h$, they both annihilate $h$; hence $Sh=0$.  Thus the order ideal $\bigcup_{c>0}[-cA,cA]$ is contained in the closed proper subspace of Hilbert--Schmidt operators that annihilate $h$, and is not dense.  So $A$ is not quasi-interior.

Conversely, assume $\ker A=\{0\}$.  Then $\Ran(A^{1/2})$ is dense in $E$.  For $z=A^{1/2}y$ one has the rank-one domination
\[
        z\otimes z \leq \|y\|^2 A
\]
as quadratic forms, since
\[
        |\langle x,z\rangle|^2
        =
        |\langle A^{1/2}x,y\rangle|^2
        \leq
        \|y\|^2\langle Ax,x\rangle.
\]
Hence every rank-one projection $z\otimes z$ with $z\in\Ran(A^{1/2})$ lies in the positive part of the principal ideal generated by $A$.  By polarization, the real linear span of these rank-one projections contains the finite-rank self-adjoint operators generated by vectors in the dense subspace $\Ran(A^{1/2})$.  Those finite-rank self-adjoint operators are dense in $\Hsa$.  Therefore the principal ideal generated by $A$ is dense, so $A$ is quasi-interior.
\end{proof}

\section*{Counterexample}

Let $D$ be the positive diagonal operator on $E=\ell_2(\N)$ given by
\[
        D e_n=\frac{1}{n^2}e_n.
\]
Then $D\in\Hsa^+$, $D$ is injective, and
\[
        \Ran(D^{1/2})
        =
        \left\{x=(x_n)\in\ell_2:\sum_{n=1}^\infty n^2 |x_n|^2<\infty\right\}
        =R.
\]
Choose, by Lemma~\ref{lem:transverse-range}, a unitary $U$ on $E$ such that $R\cap UR=\{0\}$, and define
\[
        A=D,\qquad B=UDU^*.
\]
Both $A$ and $B$ are injective positive Hilbert--Schmidt operators, hence both are quasi-interior points of $\Hsa^+$ by Lemma~\ref{lem:hs-cone}.

Now suppose $0\leq C\leq A$ and $0\leq C\leq B$ for some $C\in\Hsa^+$.  Douglas' range inclusion lemma says that for bounded operators $S,T$ on a Hilbert space, $\Ran S\subseteq\Ran T$ is equivalent to $SS^*\leq \lambda TT^*$ for some $\lambda\geq0$.  Applying this with $S=C^{1/2}$ and $T=A^{1/2}$, and then with $T=B^{1/2}$, yields
\[
        \Ran(C^{1/2})\subseteq \Ran(A^{1/2})=R,
        \qquad
        \Ran(C^{1/2})\subseteq \Ran(B^{1/2})=UR.
\]
Thus
\[
        \Ran(C^{1/2})\subseteq R\cap UR=\{0\},
\]
so $C^{1/2}=0$ and $C=0$.

This proves that the quasi-interior points $A$ and $B$ have no nonzero common positive lower bound.  In particular, they do not dominate a common quasi-interior point.

\section*{Verification and Limitations}

The cone $\Hsa^+$ is self-dual, so the source hypothesis $H^+\subseteq(H')^+$ is satisfied with equality.  The example is not a Banach lattice cone and has empty interior, which is precisely where the $L^2$ lattice infimum argument from the source fails.

The counterexample does not address the separate open problem, cited in the source from Glueck--Weber, asking whether quasi-interior and almost-interior points coincide for every generating normal cone.  It only answers the common-lower-bound question following Theorem 8.4.

\section*{Novelty Check}

Before promotion, the run indexes were searched for arXiv:2210.13566, quasi-interior points, ordered Hilbert spaces, positive Hilbert--Schmidt cones, self-dual cones, operator ranges, and Douglas range inclusion.  No exact prior packet or duplicate was found.  Web searches for close phrases such as ``quasi-interior common lower bound ordered Hilbert space'', ``positive Hilbert-Schmidt quasi-interior'', and ``operator ranges trivial intersection Hilbert space'' found related background but no explicit answer to this source question.

\section*{References}

\begin{enumerate}
    \item J. Glueck and A. Mironchenko, \emph{Stability criteria for positive semigroups on ordered Banach spaces}, arXiv:2210.13566.
    \item R. G. Douglas, \emph{On majorization, factorization, and range inclusion of operators on Hilbert space}, Proceedings of the American Mathematical Society 17 (1966), 413--415.
\end{enumerate}

\end{document}
